\documentclass[11pt,a4paper]{article}

\usepackage[utf8]{inputenc}
\usepackage[T1]{fontenc}
\usepackage[french]{babel}
\usepackage{amsmath,amssymb,amsfonts}
\usepackage{geometry}
\usepackage{graphicx}
\usepackage{booktabs}
\usepackage{hyperref}
\usepackage{array}
\usepackage{xcolor}
\usepackage{float}

\geometry{margin=2.5cm}
\graphicspath{{figures/}}

\title{
\textbf{Équations d'Einstein effectives depuis l'équilibre d'intrication}\\
\large Assemblage des limites spectrale, de Ricci et de source modulaire dans BuP
}

\author{Farid Hamdad}
\date{\today}

\begin{document}

\maketitle

\begin{abstract}
Paper~19 assemble les trois briques établies dans Papers~16--18 pour obtenir
la forme effective des équations d'Einstein dans Bottom-Up Quantum Gravity.
Paper~16 a montré que le Laplacien d'intrication reconstruit le bas spectre
du Laplacien de Laplace--Beltrami :
\[
c_NL_N\to-\Delta_g.
\]
Paper~17 a montré que la courbure d'Ollivier--Ricci contient un signal moyen
de Ricci :
\[
\frac{\kappa^{OR}}{\epsilon}
\simeq
B_N+C_N R_{\mu\nu}u^\mu u^\nu.
\]
Paper~18 a montré que la variation modulaire agit comme source effective de
la réponse de courbure :
\[
\delta W_{\rm loc}
\to
\delta S_A
\simeq
\delta\langle K_A\rangle
\to
\delta\kappa(r).
\]
En combinant ces trois flèches, Paper~19 propose l'équation effective :
\[
G_{\mu\nu}[g^{\rm ent}]
+
\Lambda_{\rm ent}g_{\mu\nu}^{\rm ent}
=
8\pi G_{\rm eff}T_{\mu\nu}^{\rm ent}
+
\mathcal{H}_{\mu\nu}.
\]
Le tenseur \(\mathcal{H}_{\mu\nu}\) encode les corrections spectrales,
topologiques, non locales et de taille finie. Dans la limite lisse et basse
énergie, \(\mathcal{H}_{\mu\nu}\to0\), et BuP admet un régime effectif de type
Einstein.
\end{abstract}

\tableofcontents

\section{Introduction}

Bottom-Up Quantum Gravity part d'une hypothèse minimale : l'espace, la
métrique et la gravité ne sont pas fondamentaux, mais émergent de la structure
d'intrication d'un état quantique global fini.

L'objet fondamental est la matrice d'information mutuelle :
\[
W_{ij}=I(i:j).
\]

À partir de cette matrice, on construit :

\[
L(W)=D-W,
\qquad
D_{ii}=\sum_j W_{ij},
\]
ainsi qu'une courbure discrète, une distance d'intrication et des termes
spectraux.

Les papiers précédents ont progressivement établi les trois ingrédients
nécessaires à une limite d'Einstein effective :

\[
L_N\to\Delta_g,
\]

\[
\kappa^{OR}\to R_{\mu\nu}u^\mu u^\nu,
\]

\[
\delta W_{\rm loc}\to T_{\mu\nu}^{\rm ent}.
\]

Paper~19 assemble ces trois résultats dans une seule structure variationnelle.

La question centrale est :
\[
\boxed{
\frac{\delta S_{\rm BuP}[W]}{\delta W_{ij}}=0
\quad
\Longrightarrow
\quad
G_{\mu\nu}
+
\Lambda_{\rm ent}g_{\mu\nu}
=
8\pi G_{\rm eff}T_{\mu\nu}^{\rm ent}
+
\mathcal{H}_{\mu\nu}
?
}
\]

\section{Action BuP discrète}

La forme minimale de l'action BuP est :
\[
S_{\rm BuP}[W]
=
\mathrm{Tr}\,L(W)^{-\beta[W]}
+
\sum_{ij}W_{ij}\kappa_{ij}[W]
+
\lambda\sum_{ij}W_{ij}d_{ij}^{2}
+
S_{\rm topo}[W].
\]

Le principe dynamique est l'équilibre :
\[
\frac{\delta S_{\rm BuP}[W]}{\delta W_{ij}}=0.
\]

Chaque terme a une fonction précise.

\subsection{Terme spectral}

Le terme
\[
\mathrm{Tr}\,L(W)^{-\beta[W]}
\]
encode la géométrie spectrale. Il contrôle la diffusion, la dimension
spectrale et la structure à bas modes du graphe d'intrication.

Paper~14 a motivé le rôle de \(\beta\) comme exposant spectral auto-contraint.
Paper~16 a ensuite montré que \(L(W)\), correctement rescalé, converge
spectralement vers le Laplacien de Laplace--Beltrami.

\subsection{Terme de courbure}

Le terme
\[
\sum_{ij}W_{ij}\kappa_{ij}[W]
\]
encode la réponse de courbure discrète. La courbure utilisée est la courbure
d'Ollivier--Ricci :
\[
\kappa_{ij}^{OR}
=
1-
\frac{W_1(m_i,m_j)}{d(i,j)}.
\]

Paper~17 a montré que cette courbure contient un signal moyen de Ricci.

\subsection{Terme de localité}

Le terme
\[
\lambda\sum_{ij}W_{ij}d_{ij}^{2}
\]
pénalise les liens trop longs et force l'émergence d'une géométrie locale. Il
joue aussi un rôle dans l'apparition d'une échelle infrarouge et donc d'un
terme cosmologique effectif.

\subsection{Terme topologique}

Le terme
\[
S_{\rm topo}[W]
\]
encode les contraintes globales de connectivité et de topologie. Il contribue
aux corrections non locales et globales dans la limite continue.

\section{Pillar I : limite spectrale}

Paper~16 a testé directement la convergence des valeurs propres de \(L_N\)
vers celles du Laplacien de Laplace--Beltrami.

La forme testée est :
\[
c_NL_N
=
c_N\frac{D_N-W_N}{\epsilon_N}
\longrightarrow
-\Delta_g.
\]

Sur le cercle \(S^1\), les 12 premiers modes non nuls sont reconstruits avec
une erreur relative moyenne :
\[
0.005072.
\]

Le premier mode donne :
\[
\lambda_1^{\rm scaled}=1.006212,
\]
à comparer à :
\[
\lambda_1^{S^1}=1.
\]

Sur le tore plat \(T^2\), les 20 premiers modes non nuls sont reconstruits
avec une erreur relative moyenne :
\[
0.026627.
\]

Le premier mode donne :
\[
\lambda_1^{\rm scaled}=40.961270,
\]
à comparer à :
\[
\lambda_1^{T^2}=4\pi^2=39.478418.
\]

Ces résultats soutiennent :
\[
L(W)\to-\Delta_g.
\]

Ainsi, le terme spectral discret peut être relié à une action géométrique
continue. En particulier, via la représentation de Mellin :
\[
\mathrm{Tr}\,L^{-\beta}
=
\frac{1}{\Gamma(\beta)}
\int_0^\infty
t^{\beta-1}
\mathrm{Tr}(e^{-tL})\,dt,
\]
et l'expansion de la trace de chaleur :
\[
\mathrm{Tr}(e^{-t\Delta_g})
\sim
(4\pi t)^{-D/2}
\sum_{n=0}^{\infty}a_n(\Delta_g)t^n,
\]
le terme spectral contient des contributions géométriques de type volume,
courbure scalaire et corrections de courbure supérieure.

\section{Pillar II : limite de Ricci}

Paper~17 a étudié la limite continue de la courbure d'Ollivier--Ricci.

La relation moyenne obtenue est :
\[
\frac{\kappa^{OR}}{\epsilon}
\simeq
B_N+C_N R_{\mu\nu}u^\mu u^\nu.
\]

Les tests v1 ont montré que la sphère présente systématiquement un signal de
courbure supérieur au tore plat :
\[
\Delta\bar\kappa_{\rm sphere-flat}=0.013383,
\]
\[
\Delta
\left\langle
\frac{\kappa}{\epsilon}
\right\rangle_{\rm sphere-flat}
=
0.197719,
\]
\[
\Delta
\left\langle
\frac{\kappa}{\ell^2}
\right\rangle_{\rm sphere-flat}
=
0.557041.
\]

La calibration affine v2 donne, à \(N=512\) :
\[
B_N=-0.301281,
\]
\[
C_N=0.391933,
\]
et :
\[
A_N=\frac{1}{C_N}=2.551455.
\]

Ainsi :
\[
\frac{\kappa^{OR}}{\epsilon}
\simeq
-0.301
+
0.392\,R_{\mu\nu}u^\mu u^\nu.
\]

Cela ne constitue pas encore une convergence point par point de la courbure,
mais cela établit un proxy moyen de Ricci.

Dans la limite continue, ce résultat permet d'identifier le secteur de
courbure :
\[
\kappa^{OR}
\to
R_{\mu\nu}u^\mu u^\nu.
\]

\section{Pillar III : source modulaire}

Paper~18 a étudié le côté source.

Le point de départ est la première loi modulaire :
\[
\delta S_A
=
\delta\langle K_A\rangle,
\]
où :
\[
K_A=-\log\rho_A.
\]

Sur une région \(A\) du graphe, Paper~18 utilise le proxy :
\[
\rho_A
=
\frac{(L_A+\mu I)^{-1}}
{\mathrm{Tr}(L_A+\mu I)^{-1}}.
\]

On définit :
\[
S_A=-\mathrm{Tr}(\rho_A\log\rho_A),
\]
et :
\[
K_A=-\log\rho_A.
\]

Après perturbation locale de \(W\), Paper~18 mesure :
\[
\delta S_A
\]
et :
\[
\delta\langle K_A\rangle
=
\mathrm{Tr}\left[(\rho'_A-\rho_A)K_A\right].
\]

Le test v1 donne, sur le tore plat :
\[
\delta S_A
\simeq
0.989481\,\delta\langle K_A\rangle,
\]
avec :
\[
R^2=0.996590.
\]

Sur la sphère :
\[
\delta S_A
\simeq
0.989718\,\delta\langle K_A\rangle,
\]
avec :
\[
R^2=0.996592.
\]

Le test v2 montre ensuite que la source modulaire prédit la réponse de
courbure proche :
\[
|\delta\langle K_A\rangle|
\to
\langle|\Delta\kappa|\rangle_{\rm near}.
\]

Sur le tore plat :
\[
R^2=0.967229,
\]
et sur la sphère :
\[
R^2=0.982252.
\]

La relation signée est quasi parfaite :
\[
r=-0.999168
\]
sur le tore plat, et :
\[
r=-0.999428
\]
sur la sphère.

Ainsi :
\[
\delta W_{\rm loc}
\to
\delta S_A
\simeq
\delta\langle K_A\rangle
\to
\delta\kappa(r).
\]

Dans la limite continue, on interprète :
\[
\delta\langle K_A\rangle
\to
\int_A \xi^\mu\delta T_{\mu\nu}^{\rm ent}\,d\Sigma^\nu.
\]

Cela définit le secteur source :
\[
T_{\mu\nu}^{\rm ent}.
\]

\section{Dictionnaire continu BuP}

Les résultats précédents permettent d'établir le dictionnaire suivant :

\[
\begin{array}{ll}
\toprule
{\rm Objet\ discret\ BuP} & {\rm Objet\ continu}\\
\midrule
W_{ij}=I(i:j) & g_{\mu\nu}^{\rm ent}\\
L(W) & -\Delta_g\\
\mathrm{Tr}\,L^{-\beta} & {\rm action\ spectrale\ gravitationnelle}\\
\kappa_{ij}^{OR} & R_{\mu\nu}u^\mu u^\nu\\
\delta\langle K_A\rangle & T_{\mu\nu}^{\rm ent}\\
\lambda\sum W_{ij}d_{ij}^2 & \Lambda_{\rm ent}g_{\mu\nu}\\
S_{\rm topo}[W] & {\rm contraintes\ topologiques}\\
{\rm corrections\ finies} & \mathcal{H}_{\mu\nu}\\
\bottomrule
\end{array}
\]

Ce dictionnaire permet de passer de :
\[
\frac{\delta S_{\rm BuP}}{\delta W_{ij}}=0
\]
à une équation tensorielle effective.

\section{Équation d'Einstein effective}

L'équation effective proposée est :
\[
\boxed{
G_{\mu\nu}[g^{\rm ent}]
+
\Lambda_{\rm ent}g_{\mu\nu}^{\rm ent}
=
8\pi G_{\rm eff}T_{\mu\nu}^{\rm ent}
+
\mathcal{H}_{\mu\nu}
}
\]

où :
\[
G_{\mu\nu}
=
R_{\mu\nu}
-
\frac{1}{2}Rg_{\mu\nu}.
\]

Le terme :
\[
\Lambda_{\rm ent}g_{\mu\nu}^{\rm ent}
\]
provient du secteur infrarouge, de la contrainte de localité et de l'énergie
de fond d'intrication.

Le terme :
\[
T_{\mu\nu}^{\rm ent}
\]
provient de la variation modulaire :
\[
\delta\langle K_A\rangle.
\]

Le tenseur :
\[
\mathcal{H}_{\mu\nu}
\]
contient les corrections qui ne se réduisent pas à la forme d'Einstein.

\section{Le tenseur de correction \(\mathcal{H}_{\mu\nu}\)}

Le terme \(\mathcal{H}_{\mu\nu}\) est nécessaire car la limite BuP n'est pas
exactement celle de la relativité générale à toutes les échelles.

Il regroupe :

\begin{enumerate}
    \item les corrections spectrales d'ordre supérieur ;
    \item les corrections non locales d'intrication ;
    \item les effets topologiques globaux ;
    \item les corrections de courbure supérieure ;
    \item les corrections de taille finie ;
    \item les effets de non-uniformité du graphe ;
    \item les écarts à la limite de variété lisse.
\end{enumerate}

Dans le régime lisse, local et basse énergie, on attend :
\[
\mathcal{H}_{\mu\nu}\to0.
\]

Alors :
\[
G_{\mu\nu}
+
\Lambda_{\rm ent}g_{\mu\nu}
=
8\pi G_{\rm eff}T_{\mu\nu}^{\rm ent}.
\]

\section{Interprétation physique}

La lecture physique de Paper~19 est la suivante.

La géométrie n'est pas posée au départ. Elle émerge du spectre de
\(L(W)\). La courbure n'est pas postulée comme un champ continu fondamental ;
elle apparaît comme limite de la courbure d'Ollivier--Ricci. La source n'est
pas introduite comme matière classique externe ; elle est reconstruite depuis
la variation modulaire d'une région du graphe.

Ainsi, la gravité apparaît comme une condition d'équilibre de l'intrication :
\[
\frac{\delta S_{\rm BuP}}{\delta W}=0.
\]

Dans cette lecture, l'équation d'Einstein effective n'est pas un postulat.
Elle est l'expression continue d'un équilibre discret entre :

\begin{itemize}
    \item spectre ;
    \item courbure ;
    \item localité ;
    \item topologie ;
    \item source modulaire.
\end{itemize}

\section{Limites}

Paper~19 est un papier d'assemblage. Il ne constitue pas encore une preuve
mathématique complète.

Les limites principales sont :

\begin{enumerate}
    \item la normalisation \(c_N\) de Paper~16 est encore ajustée numériquement ;
    \item la relation de Ricci de Paper~17 est moyenne, pas point par point ;
    \item le tenseur \(T_{\mu\nu}^{\rm ent}\) de Paper~18 est encore représenté par un proxy modulaire scalaire ;
    \item les tests sont réalisés sur des géométries contrôlées ;
    \item les graphes utilisés ne sont pas encore tous de vrais graphes d'information mutuelle issus d'états quantiques ;
    \item le tenseur de correction \(\mathcal{H}_{\mu\nu}\) reste à classifier ;
    \item \(G_{\rm eff}\) et \(\Lambda_{\rm ent}\) restent à dériver analytiquement.
\end{enumerate}

Cependant, l'assemblage est cohérent : chaque flèche nécessaire à la limite
d'Einstein possède désormais une validation numérique contrôlée.

\section{Conclusion}

Paper~19 assemble les résultats de Papers~16--18 pour proposer la forme
effective des équations d'Einstein dans BuP.

Les trois piliers sont :

\[
L_N\to\Delta_g,
\]

\[
\kappa^{OR}\to R_{\mu\nu}u^\mu u^\nu,
\]

\[
\delta W_{\rm loc}
\to
\delta\langle K_A\rangle
\to
T_{\mu\nu}^{\rm ent}.
\]

En les combinant, on obtient :
\[
G_{\mu\nu}[g^{\rm ent}]
+
\Lambda_{\rm ent}g_{\mu\nu}^{\rm ent}
=
8\pi G_{\rm eff}T_{\mu\nu}^{\rm ent}
+
\mathcal{H}_{\mu\nu}.
\]

Dans le régime lisse et basse énergie :
\[
\mathcal{H}_{\mu\nu}\to0.
\]

BuP admet donc un régime effectif de type Einstein :
\[
G_{\mu\nu}
+
\Lambda_{\rm ent}g_{\mu\nu}
=
8\pi G_{\rm eff}T_{\mu\nu}^{\rm ent}.
\]

Le résultat ne ferme pas encore la théorie. Il établit cependant que les trois
structures nécessaires à l'équation d'Einstein — opérateur géométrique,
courbure de Ricci et source énergie--impulsion — apparaissent naturellement à
partir de l'intrication.

\section*{Disponibilité du code et des données}

Les scripts et résultats de Paper~19 sont organisés dans le dossier :
\[
\texttt{papers/paper19\_effective\_einstein\_equations/}.
\]

Le script principal est :
\[
\texttt{scripts/paper19\_build\_einstein\_limit\_summary\_v1.py}.
\]

Les résultats principaux sont stockés dans :
\[
\texttt{results/einstein\_limit\_summary\_v1/}.
\]

\appendix

\section{Résumé des résultats assemblés}

\[
\begin{array}{llll}
\toprule
{\rm Paper} & {\rm Flèche} & {\rm Résultat} & {\rm Valeur}\\
\midrule
16 & L_N\to\Delta_g & {\rm erreur\ spectrale}\ S^1 & 0.005072\\
16 & L_N\to\Delta_g & {\rm erreur\ spectrale}\ T^2 & 0.026627\\
17 & \kappa^{OR}\to Ricci & B_N & -0.301281\\
17 & \kappa^{OR}\to Ricci & C_N & 0.391933\\
18 & \delta S\simeq\delta K & {\rm pente}\ T^2 & 0.989481\\
18 & \delta S\simeq\delta K & {\rm pente}\ S^2 & 0.989718\\
18 & \delta K\to\delta\kappa & R^2\ T^2 & 0.967229\\
18 & \delta K\to\delta\kappa & R^2\ S^2 & 0.982252\\
\bottomrule
\end{array}
\]

\section{Prochaines étapes}

Les prochaines étapes sont :

\begin{enumerate}
    \item dériver analytiquement \(c_N\) dans Paper~16 ;
    \item dériver analytiquement \(B_N\) et \(C_N\) dans Paper~17 ;
    \item reconstruire un vrai tenseur \(T_{\mu\nu}^{\rm ent}\) dans Paper~18 ;
    \item classifier les termes de correction \(\mathcal{H}_{\mu\nu}\) ;
    \item dériver \(G_{\rm eff}\) depuis les constantes spectrales ;
    \item relier \(\Lambda_{\rm ent}\) aux résultats cosmologiques BuP ;
    \item tester la limite effective sur des graphes d'information mutuelle quantiques ;
    \item préparer Paper~20 : corrections, prédictions et déviations observables.
\end{enumerate}

\end{document}
